\section{Introduction}\label{sec:introduction}

An Einstein metric has Ricci tensor equal to a constant multiple of the
metric. In dimension four this fixes the sum of the sectional curvatures
through each tangent direction, while leaving two independent Weyl
curvature blocks. A basic rigidity problem asks how much nonnegative
sectional curvature restricts those blocks. The product of two round
surfaces illustrates the boundary case: planes tangent to a factor have
positive curvature, whereas planes spanned by one direction from each
factor have zero curvature. We prove that, for a closed Einstein
four-manifold with positive Einstein constant, a single zero-curvature
plane already forces this product geometry on the universal cover.

\begin{theorem}[Zero-plane rigidity]\label{thm:main}
Let $(M^4,g)$ be a connected smooth closed Riemannian manifold satisfying
\[
 \operatorname{Ric}_g=3g,\qquad K_g(\sigma)\ge0
\]
for every tangent two-plane $\sigma$. If $K_g(\sigma_0)=0$ for at least
one tangent two-plane, then the universal Riemannian cover is isometric to
\[
 S^2(1/\sqrt3)\times S^2(1/\sqrt3)
\]
with its product metric. In particular, the tangent bundle of the
universal cover is the orthogonal sum of two parallel rank-two subbundles.
\end{theorem}

Here $S^2(r)$ denotes the round sphere of radius $r$, and closed means
compact without boundary. No orientation, symmetry, or topological
condition is assumed on $M$. The Einstein normalization only fixes the
scale: for $\operatorname{Ric}=\lambda g$ with $\lambda>0$, the two
radii are $1/\sqrt\lambda$.

\subsection{Context and the boundary problem}

The classification of positively curved Einstein four-manifolds has a
long history. Berger's rigidity under strict quarter pinching
\cite{Berger1961} was followed by quantitative curvature conditions of
Yang, Costa, and Cao--Tran
\cite{Yang2000,Costa2004,CaoTran2018}. Gursky--LeBrun obtained the
Fubini--Study conclusion under nonnegative sectional curvature and a
nonzero positive-definite intersection form
\cite[Theorem~A]{GurskyLeBrun1999}. Their separate Weyl-norm gaps and
analysis of the two Weyl blocks are important antecedents of the
argument here. Scalar combinations of the two Weyl norms also appear in
Yang's coupling argument~\cite[Section~6]{Yang2000} and in
Cao--Tran~\cite[Proposition~2.2]{CaoTran2018}. Under nonnegative isotropic curvature, which is a
different curvature condition, the Einstein local-symmetry theorem in
dimension four is due to Micallef--Wang
\cite[Theorem~4.4]{MicallefWang1993}; Brendle established the analogous
rigidity in higher dimensions~\cite[Theorem~1]{Brendle2010}.

Recent results treat additional topology or symmetry. Cheng proves the
oriented positive-sectional classification when the Euler characteristic
is at most three~\cite[Theorem~1.2]{Cheng2026}; Gursky--Malchiodi obtain
it under $2\chi-3|\tau|\le4$, where $\chi$ is the Euler characteristic
and $\tau$ the signature~\cite[Theorem~1.1]{GurskyMalchiodi2026}.
For nonnegative sectional curvature, Liu proves the classification of
closed simply connected manifolds with a $T^2$-invariant Einstein metric~\cite[Theorem~1.1]{Liu2025}.
Di Cerbo obtains the $S^2\times S^2$ universal cover when equality in a
Weyl-norm bound holds at every point
\cite[Proposition~3.2]{DiCerbo2026}. The hypothesis of
Theorem~\ref{thm:main} instead specifies one sectional-curvature plane
at one point.

The companion manuscript \emph{Positively curved Einstein
four-manifolds}~\cite{OpenAI2026P} proves the strictly positive
classification without additional curvature or topological restrictions.
Its coupled Weyl estimates provide the starting point of this paper.
We reproduce the tensor calculations, upper volume comparison, moment
argument, and polynomial inequalities needed here, with their
attribution at the point of use. The extra issue is their behavior at
the boundary of the sectional-curvature cone. An estimate proved for a
strictly positive metric does not by itself settle that boundary case.

The argument also gives the full classification under nonnegative
sectional curvature.

\begin{corollary}\label{cor:nonnegative-classification}
Let $(M^4,g)$ be connected, smooth, and closed, with
$\operatorname{Ric}=3g$ and nonnegative sectional curvature. Its universal
Riemannian cover is one of the round $S^4$ of curvature $1$, the complex
projective plane with its Fubini--Study metric normalized by
$\operatorname{Ric}=3g$, or $S^2(1/\sqrt3)\times S^2(1/\sqrt3)$.
\end{corollary}
Its proof is given at the end of Section~\ref{sec:classification}.

\subsection{Reduction and proof strategy}

The universal Riemannian cover is complete and has $\operatorname{Ric}=3g$,
so Bonnet--Myers makes it compact~\cite{Myers1941}. It is simply connected
and therefore orientable. We work on this cover from now on and keep the
notation $(M,g)$. The zero-curvature plane lifts to it.

Choose an orientation and let $\Lambda^+$ and $\Lambda^-$ be the
rank-three eigenspaces of the Hodge star on two-forms. The Einstein
curvature operator has blocks $\operatorname{Id}+T$ and
$\operatorname{Id}+U$, where $T,U$ are symmetric and trace-free.
They are the self-dual and anti-self-dual Weyl blocks. Set
\[
 v=\frac{|T|}{\sqrt6},\qquad b=\frac{|U|}{\sqrt6},\qquad
 x=-\lambda_{\min}(T),\qquad y=-\lambda_{\min}(U).
\]
All matrix norms are Hilbert--Schmidt norms. As proved in
Lemma~\ref{lem:curvature-blocks}, the smallest sectional curvature is
$1-(x+y)/2$, and
\[
 v\le x\le2v,\qquad b\le y\le2b,\qquad x+y\le2.
\]
The stronger spectral inequality $x+y\le2$, not merely $v+b\le2$,
is retained in the tensor estimates.

If one Weyl block vanished identically, the metric would be
half-conformally flat. The classical compact positive-scalar Einstein
classification of Hitchin~\cite{Hitchin1981}, with conformal Einstein
rigidity~\cite{Obata1962}, would identify our simply connected cover
with the standard $S^4$ or $\mathbb{CP}^2$ up to scale; see also
\cite[Remark~7]{Derdzinski1983}. Both have strictly positive sectional
curvature. We may therefore assume throughout the remaining proof that
neither block vanishes identically. The zero-plane hypothesis is used
only in this reduction: the remaining estimates assume only nonnegative
sectional curvature and that neither block vanishes identically.
Individual blocks may still vanish at some points.

Write $\mathbf E f$ for the integral of $f$ divided by the volume of $M$.
The proof now has three stages.
First, the separate gaps $\mathbf E v^2,\mathbf E b^2\ge1$,
a weighted variance inequality, and an upper volume deficit force
\[
 \mathbf E v^2=\mathbf E b^2.
\]
Second, a weighted Hessian identity and explicit polynomial inequalities
couple the two blocks and force $dv=db=0$ almost everywhere.
Third, the gaps and $v+b\le2$ give $v=b=1$. The weak norm equations
then imply that both Weyl blocks are parallel, each with spectrum
$(2,-1,-1)$. Their simple eigenlines produce the parallel tangent
splitting.

Three features allow these estimates to reach the boundary. The upper
volume estimate integrates over the domain before the first conjugate
point and needs no lower bound on the injectivity radius. The weighted
Hessian identity needs a nonnegative quadratic form, so it remains valid
when that form has a kernel. At a zero of one Weyl block, the pointwise
algebraic estimates are obtained by introducing a small nonzero block
and shrinking the opposite block while keeping its derivative data fixed.
The decisive strictness then comes from the fact that the gradient of a
Lipschitz function vanishes almost everywhere on a level set. This forces
constant Weyl norms without requiring strict sectional positivity.

Section~\ref{sec:preliminaries} develops the curvature identities,
including the weak norm equations at zeros. Sections~\ref{sec:analytic}
and~\ref{sec:volume} establish the scalar estimates and moment balance.
Sections~\ref{sec:hessian} and~\ref{sec:coupling} give the Hessian and
coupling inequalities. Section~\ref{sec:classification} resolves their
equality case and proves Theorem~\ref{thm:main}.
Appendix~\ref{sec:certificates} supplies the exact polynomial
certificates, including all boundary faces.
